% ============================================================
% AUTO-GENERATED by rm3_ablation_and_integration.py
% DO NOT EDIT MANUALLY — re-run the script to refresh.
% ============================================================

% Table IV.13 — RM3 summary across all four blocks
\begin{table}[!htbp]
\centering
\caption{Ringkasan evaluasi RM3 pada penulisan graf, pembaruan
\textit{lifecycle}, retrieval, dan penggunaan memori pada jawaban}
\label{tab:v2-memory-result}

{\footnotesize
\renewcommand{\arraystretch}{1.10}
\setlength{\tabcolsep}{4pt}

\begin{tabularx}{\textwidth}{
|>{\raggedright\arraybackslash}X
|>{\centering\arraybackslash}p{2.40cm}
|>{\centering\arraybackslash}p{1.70cm}|
}
\hline

\rowcolor{tableheadgray}
\textbf{Metrik}
&
\textbf{Kondisi}
&
\textbf{Nilai}
\\
\hline

\multicolumn{3}{|l|}{
\textit{Penulisan graf}, $n = 12$ skenario
}
\\
\hline

F1 node
&
Tidak berlaku
&
0,620
\\
\hline

Presisi node
&
Tidak berlaku
&
0,556
\\
\hline

Recall node
&
Tidak berlaku
&
0,702
\\
\hline

F1 relasi
&
Tidak berlaku
&
0,361
\\
\hline

Akurasi label \texttt{subject\_role}
&
Tidak berlaku
&
85,7\%
\\
\hline

Akurasi label \texttt{trigger\_category}
&
Tidak berlaku
&
100,0\%
\\
\hline

\multicolumn{3}{|l|}{
\textit{Pembaruan lifecycle}, $n = 12$ skenario
}
\\
\hline

Ketepatan tindakan pembaruan
&
Tidak berlaku
&
33,3\%
\\
\hline

Ketepatan konten baru
&
Tidak berlaku
&
50,0\%
\\
\hline

Konten lama berhasil dinonaktifkan
&
Tidak berlaku
&
25,0\%
\\
\hline

\multicolumn{3}{|l|}{
\textit{Retrieval}, $n = 42$ kueri pada dua kondisi
}
\\
\hline

P@5
&
Vektor saja
&
0,000
\\
\hline

P@5
&
Hibrid
&
0,052
\\
\hline

Recall@5
&
Vektor saja
&
0,286
\\
\hline

Recall@5
&
Hibrid
&
0,357
\\
\hline

MRR
&
Vektor saja
&
0,000
\\
\hline

MRR
&
Hibrid
&
0,190
\\
\hline

nDCG@5
&
Vektor saja
&
0,286
\\
\hline

nDCG@5
&
Hibrid
&
0,469
\\
\hline

Tingkat memori usang
&
Vektor saja
&
0,0\%
\\
\hline

Tingkat memori usang
&
Hibrid
&
0,0\%
\\
\hline

Keberhasilan kueri relasional atau kausal
&
Vektor saja
&
0,0\%
\\
\hline

Keberhasilan kueri relasional atau kausal
&
Hibrid
&
29,2\%
\\
\hline

Keberhasilan abstensi
&
Vektor saja
&
28,6\%
\\
\hline

Keberhasilan abstensi
&
Hibrid
&
45,2\%
\\
\hline

\multicolumn{3}{|l|}{
\textit{Penggunaan memori pada jawaban},
$n = 24$ kasus pada dua kondisi
}
\\
\hline

Jawaban berdasar memori
&
Vektor saja
&
95,8\%
\\
\hline

Jawaban berdasar memori
&
Hibrid
&
95,8\%
\\
\hline

Tingkat kontradiksi
&
Vektor saja
&
4,2\%
\\
\hline

Tingkat kontradiksi
&
Hibrid
&
8,3\%
\\
\hline

Keberhasilan relasional atau kausal
&
Vektor saja
&
62,5\%
\\
\hline

Keberhasilan relasional atau kausal
&
Hibrid
&
66,7\%
\\
\hline

Personalisasi tanpa dasar
&
Vektor saja
&
29,2\%
\\
\hline

Personalisasi tanpa dasar
&
Hibrid
&
29,2\%
\\
\hline

Tingkat keyakinan usang
&
Vektor saja
&
8,3\%
\\
\hline

Tingkat keyakinan usang
&
Hibrid
&
8,3\%
\\
\hline

Abstensi benar
&
Vektor saja
&
25,0\%
\\
\hline

Abstensi benar
&
Hibrid
&
25,0\%
\\
\hline

\end{tabularx}

\vspace{2pt}

\begin{minipage}{\textwidth}
\footnotesize
\textit{Catatan:}
Kondisi tidak berlaku digunakan pada evaluasi yang tidak membandingkan
retrieval hibrid dengan retrieval vektor saja. Nilai yang lebih tinggi
menunjukkan hasil yang lebih baik, kecuali pada tingkat kontradiksi,
tingkat memori usang, personalisasi tanpa dasar, dan tingkat keyakinan
usang.
\end{minipage}
}
\end{table}
% Table IV.14 — Ablation: hybrid vs vector-only
\begin{table}[!htbp]
\centering
\caption{Hasil ablasi menggunakan \textit{paired bootstrap} dan
\textit{sign test}: kondisi hibrid dibandingkan kondisi vektor saja}
\label{tab:v2-ablation-result}

{\footnotesize
\renewcommand{\arraystretch}{1.15}

\begin{tabularx}{\textwidth}{
|>{\raggedright\arraybackslash}p{1.35cm}
|>{\raggedright\arraybackslash}p{2.65cm}
|>{\centering\arraybackslash}p{1.30cm}
|>{\centering\arraybackslash}p{1.30cm}
|>{\centering\arraybackslash}p{1.10cm}
|>{\raggedright\arraybackslash}X|
}
\hline

\rowcolor{tableheadgray}
\textbf{Blok}
&
\textbf{Metrik}
&
\textbf{Vektor saja}
&
\textbf{Hibrid}
&
\textbf{$\Delta$}
&
\textbf{Uji statistik}
\\
\hline

\multicolumn{6}{|l|}{
\textit{Retrieval}: \textit{paired bootstrap}, $n = 42$ kueri
}
\\
\hline

Retrieval
&
P@5
&
0,000
&
0,052
&
0,052
&
95\% CI: [0,024; 0,086]
\\
\hline

Retrieval
&
Recall@5
&
0,286
&
0,357
&
0,071
&
95\% CI: [$-0,095$; 0,238]
\\
\hline

Retrieval
&
MRR
&
0,000
&
0,190
&
0,190
&
95\% CI: [0,083; 0,310]
\\
\hline

Retrieval
&
nDCG@5
&
0,286
&
0,469
&
0,184
&
95\% CI: [0,082; 0,295]
\\
\hline

\multicolumn{6}{|l|}{
\textit{Retrieval}: \textit{exact sign test} untuk metrik biner
}
\\
\hline

Retrieval
&
Relasional atau kausal
&
0,0\%
&
29,2\%
&
0,292
&
$p = 0,02$
\\
\hline

Retrieval
&
Keberhasilan abstensi
&
28,6\%
&
45,2\%
&
0,166
&
$p = 0,14$
\\
\hline

Retrieval
&
Kejadian memori usang
&
0,0\%
&
0,0\%
&
0,000
&
$p = 1,00$
\\
\hline

\multicolumn{6}{|l|}{
\textit{Penggunaan memori pada jawaban}: \textit{exact sign test},
$n = 24$ kasus
}
\\
\hline

Jawaban
&
Jawaban berdasar memori
&
95,8\%
&
95,8\%
&
0,000
&
$p = 1,00$
\\
\hline

Jawaban
&
Kontradiksi
&
4,2\%
&
8,3\%
&
0,041
&
$p = 1,00$
\\
\hline

Jawaban
&
Relasional atau kausal
&
62,5\%
&
66,7\%
&
0,042
&
$p = 1,00$
\\
\hline

Jawaban
&
Personalisasi tanpa dasar
&
29,2\%
&
29,2\%
&
0,000
&
$p = 1,00$
\\
\hline

Jawaban
&
Keyakinan usang
&
8,3\%
&
8,3\%
&
0,000
&
$p = 1,00$
\\
\hline

Jawaban
&
Abstensi benar
&
25,0\%
&
25,0\%
&
0,000
&
$p = 1,00$
\\
\hline

\end{tabularx}

\vspace{2pt}

\begin{minipage}{\textwidth}
\footnotesize
\textit{Catatan:}
Nilai $\Delta$ dihitung sebagai nilai kondisi hibrid dikurangi nilai
kondisi vektor saja. Untuk metrik kesalahan, seperti kontradiksi,
personalisasi tanpa dasar, kejadian memori usang, dan keyakinan usang,
nilai $\Delta$ positif menunjukkan peningkatan tingkat kesalahan.
\end{minipage}
}
\end{table}
% Table IV.15 — Inter-judge Cohen's kappa (answer-usage)
\begin{table}[!htbp]
\centering
\caption{Kesepakatan antarpenilai pada evaluasi penggunaan memori
dalam jawaban}
\label{tab:rm3-answer-kappa}

{\footnotesize
\renewcommand{\arraystretch}{1.15}

\begin{tabularx}{0.78\textwidth}{
|>{\raggedright\arraybackslash}X
|>{\centering\arraybackslash}p{3.10cm}|
}
\hline

\rowcolor{tableheadgray}
\textbf{Label penilaian}
&
\textbf{Cohen's $\kappa$}
\\
\hline

Jawaban berdasar memori
&
0,048
\\
\hline

Kontradiksi
&
0,000
\\
\hline

Keberhasilan relasional atau kausal
&
0,483
\\
\hline

Personalisasi tanpa dasar
&
0,628
\\
\hline

Abstensi benar
&
1,000
\\
\hline

Keyakinan usang
&
0,000
\\
\hline

\rowcolor{tableheadgray}
\textbf{Rerata deskriptif antarlabel}
&
\textbf{0,360}
\\
\hline

\end{tabularx}

\vspace{2pt}

\begin{minipage}{0.78\textwidth}
\footnotesize
\textit{Catatan:}
Nilai pada baris terakhir merupakan rerata deskriptif dari enam nilai
Cohen's $\kappa$, bukan nilai kappa gabungan yang dihitung dari seluruh
label sebagai satu variabel.
\end{minipage}
}
\end{table}